A comparação entre detectores one-class para batimentos ECG anômalos exige
avaliar, no mesmo protocolo, como a escolha do algoritmo responde a
representações progressivamente mais ricas do sinal e a diferentes custos de
inferência. A análise usa a MIT-BIH Arrhythmia Database em formato WFDB, com
canal MLII, janelas centradas nas anotações de batimento, separação
inter-paciente e treinamento restrito a batimentos normais. O protocolo rotula
batimentos \texttt{N} como normais e demais símbolos válidos não-paced como
anomalias apenas na etapa de avaliação. Quatro estágios cumulativos organizam a
engenharia de features, partindo de descritores morfológicos e intervalo RR no
sinal bruto, seguindo para o sinal filtrado por FIR, depois para descritores
espectrais e STFT, e por fim para respostas de filtros de Gabor 1D em janelas
aproximadas P, QRS e T. One-Class SVM, Isolation Forest e Local Outlier Factor
em modo novelty são comparados com foco em métricas adequadas ao
desbalanceamento, sobretudo PR-AUC, além de ROC-AUC, precision, recall,
F1-score, matriz de confusão e tempo de escore por batimento. Os resultados
preliminares indicam comportamento não monotônico entre algoritmos e estágios
de representação, com variações de desempenho que dependem da interação entre
detector e conjunto de features. A comparação fornece uma base reprodutível
para discutir quais abordagens one-class são mais competitivas nesse cenário e
qual par representação-detector oferece melhor compromisso entre separação de
anomalias e custo computacional.
